\documentclass[11pt]{article}
\usepackage[margin=1in]{geometry}
\usepackage{amsmath,amssymb}
\usepackage{bm}
\usepackage{array}
\usepackage{booktabs}
\usepackage[hidelinks]{hyperref}
\title{Part 2: Detailed Inverse PINN-RC Formulation}
\author{}
\date{}
\begin{document}
\maketitle

\section*{Part 2 --- Detailed Inverse PINN-RC Formulation}

We now define \textbf{only the inverse PINN-RC baseline}, using the Part-1 RC architectures without changing any of them.

The role of the inverse PINN is very specific:
\begin{align}
\boxed{\text{use a neural trajectory surrogate only to identify the physical RC parameters}}
\end{align}
and, after identification,
\begin{align}
\boxed{\text{discard the PINN trajectory network and retain only the identified RC model.}}
\end{align}
The inverse PINN-RC is the interpretable physics baseline, and the deployed model is the identified RC ODE rather than the PINN itself.

\section*{1. What the inverse PINN is actually solving}
For a given RC architecture, let the true physical state be
\begin{align}
x(t)\in\mathbb R^{n_x},
\end{align}
the measured output be
\begin{align}
y(t)=Hx(t),
\end{align}
and the measured forcing vector be
\begin{align}
v(t).
\end{align}
The physical RC model from Part 1 is written generally as
\begin{align}
\boxed{\mathcal F\left(x,\dot x,v;\theta\right)=0}
\tag{P2.1}
\end{align}
where \(\theta\) contains the unknown RC parameters. For example,
\begin{align}
\theta&=\{R,C,\eta_r\},
\end{align}
or
\begin{align}
\theta&=\{R_{Do},R_{Ko},R_{DK},C_D,C_K,\eta_{r,D},\eta_{r,K}\},
\end{align}
or the corresponding 2R2C parameter set. For DEP2, \(\theta\) also contains \(\alpha_c,\alpha_r\).

The PINN introduces a neural approximation
\begin{align}
\boxed{\widehat x_\phi(t)}
\tag{P2.2}
\end{align}
to the unknown physical state trajectory. The neural parameters \(\phi\) and physical parameters \(\theta\) are estimated jointly:
\begin{align}
\boxed{(\widehat\phi,\widehat\theta)=\arg\min_{\phi,\theta}\mathcal L_{\mathrm{InvPINN}}(\phi,\theta).}
\tag{P2.3}
\end{align}
But only
\begin{align}
\boxed{\widehat\theta}
\end{align}
is retained afterward.

\section*{2. Phase-D training data used by the inverse PINN}
Let
\begin{align}
t_k=t_0+k\Delta t,
\end{align}
with
\begin{align}
\boxed{\Delta t=300\ {\rm s}.}
\tag{P2.4}
\end{align}
The Phase-D \texttt{ML\_SciML/l1\_h1} silo provides at each sample the measured output
\begin{align}
y_k=y(t_k),
\end{align}
the measured forcing
\begin{align}
v_k=v(t_k),
\end{align}
and the one-step target
\begin{align}
y_{k+1}.
\end{align}
For PINN parameter identification we use the \texttt{lag\_0} observed state and forcing values. Define
\begin{align}
\boxed{\mathcal I_{\rm tr}=\left\{k:\operatorname{partition}(k)=\text{train}\right\}.}
\tag{P2.5}
\end{align}
No validation or test temperature observation enters the inverse-PINN training objective. The \texttt{target\_1} columns are therefore not required to define the PINN physics residual. They remain useful for validation, test evaluation, Sim1, and downstream comparison.

\section*{3. Why the PINN trajectory network depends on time only}
For a fixed RestaurantFastFood/Buffalo realization, the measured exogenous forcing trajectories are known functions of time. Therefore we define
\begin{align}
\boxed{\widehat x_\phi(t)=\mathcal N_\phi(t)}
\tag{P2.6}
\end{align}
rather than
\begin{align}
\widehat x_\phi=\mathcal N_\phi(x,u,d).
\end{align}
The inverse PINN is not learning the dynamical vector field. It is learning a differentiable representation of the realized trajectory so that
\begin{align}
\frac{d\widehat x_\phi}{dt}
\end{align}
can be obtained by automatic differentiation. The learned physical parameters then make that trajectory satisfy the RC equations.

\section*{4. Time normalization and exact derivative conversion}
Define
\begin{align}
t_{\min}&=\min_{k\in\mathcal I_{\rm tr}}t_k,\\
t_{\max}&=\max_{k\in\mathcal I_{\rm tr}}t_k,\\
t_c&=\frac{t_{\max}+t_{\min}}{2},\\
s_t&=\frac{t_{\max}-t_{\min}}{2}.
\end{align}
Define normalized time
\begin{align}
\boxed{\tau=\frac{t-t_c}{s_t}.}
\tag{P2.7}
\end{align}
Hence
\begin{align}
\tau\in[-1,1].
\end{align}
By the chain rule,
\begin{align}
\boxed{\frac{d\widehat x_\phi}{dt}=\frac{1}{s_t}\frac{\partial\widehat x_\phi}{\partial\tau}.}
\tag{P2.8}
\end{align}
This factor must be included explicitly; otherwise the learned capacitances would absorb the time-normalization error.

\section*{5. State-output normalization without corrupting physics}
For numerical conditioning, let \(\mu_x\) and diagonal scale matrix \(S_x\) be computed from the training partition only. Let the raw network output be \(z_\phi(\tau)\). We define the physical state represented by the PINN as
\begin{align}
\boxed{\widehat x_\phi(t)=\mu_x+S_xz_\phi(\tau).}
\tag{P2.9}
\end{align}
Therefore
\begin{align}
\boxed{\dot{\widehat x}_\phi(t)=\frac{S_x}{s_t}\frac{\partial z_\phi}{\partial\tau}.}
\tag{P2.10}
\end{align}
Every physical residual is evaluated using \(\widehat x_\phi\) in degrees Celsius, \(\dot{\widehat x}_\phi\) in degrees Celsius per second, \(Q\) in watts, \(R\) in kelvin per watt, and \(C\) in joules per kelvin. We do not write the RC equations directly in arbitrary normalized coordinates.

\section*{6. Detailed derivation of the output scaling matrix \(S_y\)}
The data loss must combine measured temperature channels without allowing one output to dominate merely because its numerical spread is larger. Let the number of measured outputs be
\begin{align}
n_y.
\end{align}
At each training time \(t_k\),
\begin{align}
y_k=\begin{bmatrix}y_{k,1}\\ \vdots\\ y_{k,n_y}\end{bmatrix}.
\end{align}
All normalization quantities are computed only from \(\mathcal I_{\rm tr}\). Define
\begin{align}
N_{\rm tr}=|\mathcal I_{\rm tr}|.
\end{align}
For each measured output channel \(j\), define the training mean
\begin{align}
\boxed{\mu_{y,j}=\frac{1}{N_{\rm tr}}\sum_{k\in\mathcal I_{\rm tr}}y_{k,j}.}
\tag{P2.11}
\end{align}
The centered training observations are
\begin{align}
\widetilde y^{\,c}_{k,j}=y_{k,j}-\mu_{y,j}.
\tag{P2.12}
\end{align}
Define the training variance as
\begin{align}
\boxed{\sigma_{y,j}^2=\frac{1}{N_{\rm tr}}\sum_{k\in\mathcal I_{\rm tr}}\left(y_{k,j}-\mu_{y,j}\right)^2.}
\tag{P2.13}
\end{align}
The training standard deviation is
\begin{align}
\boxed{\sigma_{y,j}=\sqrt{\frac{1}{N_{\rm tr}}\sum_{k\in\mathcal I_{\rm tr}}\left(y_{k,j}-\mu_{y,j}\right)^2}.}
\tag{P2.14}
\end{align}
To avoid numerical singularity if a channel has extremely small variance, introduce
\begin{align}
\epsilon_y>0.
\end{align}
The scale used in the loss is
\begin{align}
\boxed{s_{y,j}=\max\left(\sigma_{y,j},\epsilon_y\right).}
\tag{P2.15}
\end{align}
The output scaling matrix is therefore
\begin{align}
\boxed{S_y=\operatorname{diag}\left(s_{y,1},s_{y,2},\ldots,s_{y,n_y}\right).}
\tag{P2.16}
\end{align}
Its inverse is
\begin{align}
\boxed{S_y^{-1}=\operatorname{diag}\left(\frac{1}{s_{y,1}},\frac{1}{s_{y,2}},\ldots,\frac{1}{s_{y,n_y}}\right).}
\tag{P2.17}
\end{align}
For prediction error
\begin{align}
e_k=\widehat y_k-y_k,
\end{align}
the scaled error is
\begin{align}
\boxed{\widetilde e_k=S_y^{-1}e_k.}
\tag{P2.18}
\end{align}
Thus
\begin{align}
\boxed{\left\|S_y^{-1}(\widehat y_k-y_k)\right\|_2^2=\sum_{j=1}^{n_y}\left(\frac{\widehat y_{k,j}-y_{k,j}}{s_{y,j}}\right)^2.}
\tag{P2.19}
\end{align}
Each output therefore contributes in units of its own training standard deviation.

For a single-output architecture such as all-to-one or an IND zone,
\begin{align}
n_y=1,
\end{align}
and
\begin{align}
\boxed{S_y=\begin{bmatrix}s_T\end{bmatrix}.}
\tag{P2.20}
\end{align}
The scaled squared error becomes
\begin{align}
\boxed{\left\|S_y^{-1}(\widehat y_k-y_k)\right\|_2^2=\left(\frac{\widehat T_k-T_k}{s_T}\right)^2.}
\tag{P2.21}
\end{align}
For coupled DEP1/DEP2 with Dining and Kitchen outputs,
\begin{align}
n_y=2,
\end{align}
and
\begin{align}
\boxed{S_y=\begin{bmatrix}s_{T_D}&0\\0&s_{T_K}\end{bmatrix}.}
\tag{P2.22}
\end{align}
Hence
\begin{align}
\boxed{\begin{aligned}
\left\|S_y^{-1}(\widehat y_k-y_k)\right\|_2^2={}&\left(\frac{\widehat T_{D,k}-T_{D,k}}{s_{T_D}}\right)^2\\
&+\left(\frac{\widehat T_{K,k}-T_{K,k}}{s_{T_K}}\right)^2.
\end{aligned}}
\tag{P2.23}
\end{align}
The mean vector is
\begin{align}
\boxed{\mu_y=\begin{bmatrix}\mu_{y,1}\\ \vdots\\ \mu_{y,n_y}\end{bmatrix}.}
\tag{P2.24}
\end{align}
If explicit output standardization is used, define
\begin{align}
\boxed{\widetilde y_k=S_y^{-1}(y_k-\mu_y).}
\tag{P2.25}
\end{align}
The corresponding inverse transformation is
\begin{align}
\boxed{y_k=\mu_y+S_y\widetilde y_k.}
\tag{P2.26}
\end{align}
The means cancel from the prediction-error normalization because
\begin{align}
\left[(\widehat y_k-\mu_y)-(y_k-\mu_y)\right]=\widehat y_k-y_k.
\tag{P2.27}
\end{align}
Therefore the data loss needs \(S_y\) explicitly, while \(\mu_y\) is needed only if the network inputs/outputs themselves are standardized.

The leakage rule is
\begin{align}
\boxed{\mu_y,\;S_y\text{ are computed from training data only.}}
\tag{P2.28}
\end{align}
Validation and test observations are transformed using the frozen training values. We do not recompute validation/test means or scales. The final output-scaling contract is
\begin{align}
\boxed{S_y=\operatorname{diag}\left[\max(\sigma_{y,1},\epsilon_y),\ldots,\max(\sigma_{y,n_y},\epsilon_y)\right].}
\tag{P2.29}
\end{align}

\section*{7. Observed and hidden PINN outputs}
For 1R1C, all physical states are observed. For all-to-one,
\begin{align}
\boxed{\widehat x_\phi=\widehat T_A.}
\tag{P2.30}
\end{align}
For coupled identity,
\begin{align}
\boxed{\widehat x_\phi=\begin{bmatrix}\widehat T_D\\\widehat T_K\end{bmatrix}.}
\tag{P2.31}
\end{align}
No hidden PINN state exists.

For 2R2C, the PINN produces both observed air temperature and hidden mass temperature. For all-to-one,
\begin{align}
\boxed{\widehat x_\phi=\begin{bmatrix}\widehat T_{a,A}\\\widehat T_{m,A}\end{bmatrix}.}
\tag{P2.32}
\end{align}
For coupled identity,
\begin{align}
\boxed{\widehat x_\phi=\begin{bmatrix}\widehat T_{a,D}\\\widehat T_{m,D}\\\widehat T_{a,K}\\\widehat T_{m,K}\end{bmatrix}.}
\tag{P2.33}
\end{align}
The observation matrix remains
\begin{align}
\boxed{H=\begin{bmatrix}1&0&0&0\\0&0&1&0\end{bmatrix}.}
\tag{P2.34}
\end{align}
Thus
\begin{align}
\boxed{\widehat y_\phi=H\widehat x_\phi.}
\tag{P2.35}
\end{align}
The mass temperatures are not given artificial targets; they are inferred through the physical residual.

\section*{8. Physical parameter transformations}
For every resistance,
\begin{align}
\boxed{R_j=R_{j,\min}+\operatorname{softplus}(\rho_{R_j}).}
\tag{P2.36}
\end{align}
For every capacitance,
\begin{align}
\boxed{C_j=C_{j,\min}+\operatorname{softplus}(\rho_{C_j}).}
\tag{P2.37}
\end{align}
with
\begin{align}
R_{j,\min}>0,\qquad C_{j,\min}>0.
\end{align}
The optimization variables are therefore
\begin{align}
\rho_R,\rho_C\in\mathbb R.
\end{align}

\section*{9. Radiative parameter \(\eta_r\)}
For 1R1C,
\begin{align}
\eta_r^{(1C)}=\begin{cases}
1,&\texttt{full},\\
0,&\texttt{zero},\\
\eta_{\rm fixed},&\texttt{fixed},\\
\sigma(\rho_\eta),&\texttt{learnable}.
\end{cases}
\tag{P2.38}
\end{align}
For 2R2C,
\begin{align}
\eta_r^{(2C)}=\begin{cases}
1,&\texttt{mass\_only},\\
0,&\texttt{air\_only},\\
\eta_{\rm fixed},&\texttt{fixed},\\
\sigma(\rho_\eta),&\texttt{learnable}.
\end{cases}
\tag{P2.39}
\end{align}
If the option is not \texttt{learnable}, \(\rho_\eta\) does not belong to the optimization variables.

\section*{10. DEP2 allocation parameters}
DEP2 uses
\begin{align}
\boxed{\lambda_{c,D}=2\sigma(\alpha_c),}
\tag{P2.40}
\end{align}
\begin{align}
\boxed{\lambda_{c,K}=2[1-\sigma(\alpha_c)],}
\tag{P2.41}
\end{align}
\begin{align}
\boxed{\lambda_{r,D}=2\sigma(\alpha_r),}
\tag{P2.42}
\end{align}
\begin{align}
\boxed{\lambda_{r,K}=2[1-\sigma(\alpha_r)].}
\tag{P2.43}
\end{align}
Hence
\begin{align}
\boxed{\lambda_{c,D}+\lambda_{c,K}=2,}
\tag{P2.44}
\end{align}
and
\begin{align}
\boxed{\lambda_{r,D}+\lambda_{r,K}=2.}
\tag{P2.45}
\end{align}
The unconstrained trainable variables are
\begin{align}
\boxed{\alpha_c,\alpha_r\in\mathbb R.}
\tag{P2.46}
\end{align}

\section*{11. Important identifiability safeguard for DEP2}
For DEP2 1R1C, the radiative terms contain the products
\begin{align}
\eta_{r,D}\lambda_{r,D}\bar Q_r,
\end{align}
and
\begin{align}
\eta_{r,K}\lambda_{r,K}\bar Q_r.
\end{align}
If both \(\eta_r\) and \(\lambda_r\) are freely learned, the temperature equations primarily see their products. Therefore, for the primary DEP2 inverse-PINN experiment,
\begin{align}
\boxed{\text{learn }\lambda_r\quad\Longrightarrow\quad\eta_r^{(1C)}=1.}
\tag{P2.47}
\end{align}
Thus
\begin{align}
\boxed{\eta_{r,D}=\eta_{r,K}=1}
\tag{P2.48}
\end{align}
and \(\alpha_r\) is estimated. An optional ablation may instead fix
\begin{align}
\lambda_{r,D}=\lambda_{r,K}=1
\end{align}
and learn \(\eta_r\), but both are not learned simultaneously in the primary DEP2 1R1C model.

\section*{12. Data loss}
Define
\begin{align}
\widehat y_k=H\widehat x_\phi(t_k).
\end{align}
Using the training-only matrix \(S_y\) derived above,
\begin{align}
\boxed{\mathcal L_y=\frac{1}{|\mathcal I_{\rm tr}|}\sum_{k\in\mathcal I_{\rm tr}}\left\|S_y^{-1}(\widehat y_k-y_k)\right\|_2^2.}
\tag{P2.49}
\end{align}
For all-to-one \(n_y=1\); for DEP1/DEP2 \(n_y=2\); for IND, Dining and Kitchen have separate scalar losses.

\section*{13. Physics-residual normalization}
All RC residuals have units of power. For each modeled zone \(i\), define the training-only characteristic heat scale
\begin{align}
\boxed{q_i^\star=\sqrt{\frac{1}{N_{\rm tr}}\sum_{k\in\mathcal I_{\rm tr}}\left(Q_{c,i,k}^2+Q_{r,i,k}^2\right)}+\epsilon_q.}
\tag{P2.50}
\end{align}
For DEP2 use
\begin{align}
\boxed{q_i^\star=\sqrt{\frac{1}{N_{\rm tr}}\sum_{k\in\mathcal I_{\rm tr}}\left(Q_{AC,i,k}^2+(\bar Q_{c,k}^{nh})^2+(\bar Q_{r,k})^2\right)}+\epsilon_q.}
\tag{P2.51}
\end{align}
The same zone scale is used for both air and mass residuals in 2R2C.

\section*{14. All-to-One 1R1C inverse PINN}
The network is
\begin{align}
\boxed{\widehat T_A(t)=\mathcal N_\phi(t).}
\tag{P2.52}
\end{align}
The derivative is
\begin{align}
\boxed{\dot{\widehat T}_A=\frac{d\widehat T_A}{dt}.}
\tag{P2.53}
\end{align}
Define
\begin{align}
Q_{c,A}&=Q_{ZIC,A}+Q_{Sol1,A}+Q_{AC,A},\\
Q_{r,A}&=Q_{ZIR,A}+Q_{Sol2,A}.
\end{align}
The residual is
\begin{align}
\boxed{\begin{aligned}
r_A(t)={}&C_A\dot{\widehat T}_A-\frac{T_o-\widehat T_A}{R_{Ao}}-Q_{c,A}-\eta_{r,A}Q_{r,A}.
\end{aligned}}
\tag{P2.54}
\end{align}
The physics loss is
\begin{align}
\boxed{\mathcal L_f=\frac{1}{N_{\rm tr}}\sum_{k\in\mathcal I_{\rm tr}}\left(\frac{r_A(t_k)}{q_A^\star}\right)^2.}
\tag{P2.55}
\end{align}
The inverse problem is
\begin{align}
\boxed{\min_{\phi,\rho_{R_{Ao}},\rho_{C_A},\rho_{\eta_A}}\lambda_y\mathcal L_y+\lambda_f\mathcal L_f,}
\tag{P2.56}
\end{align}
where \(\rho_{\eta_A}\) is included only if \texttt{eta\_rad=learnable}. The retained model is
\begin{align}
\boxed{\dot T_A=\frac{1}{\widehat C_A}\left[\frac{T_o-T_A}{\widehat R_{Ao}}+Q_{c,A}+\widehat\eta_{r,A}Q_{r,A}\right].}
\tag{P2.57}
\end{align}
The neural trajectory network is discarded.

\section*{15. All-to-One 2R2C inverse PINN}
The network has two outputs:
\begin{align}
\boxed{\widehat x_A(t)=\begin{bmatrix}\widehat T_{a,A}(t)\\\widehat T_{m,A}(t)\end{bmatrix}.}
\tag{P2.58}
\end{align}
Its derivative is
\begin{align}
\boxed{\dot{\widehat x}_A=\begin{bmatrix}\dot{\widehat T}_{a,A}\\\dot{\widehat T}_{m,A}\end{bmatrix}.}
\tag{P2.59}
\end{align}
Only the air temperature has a measurement loss:
\begin{align}
\boxed{\mathcal L_y=\frac{1}{N_{\rm tr}}\sum_{k\in\mathcal I_{\rm tr}}\left(\frac{\widehat T_{a,A}(t_k)-T_{A,k}}{s_{T_A}}\right)^2.}
\tag{P2.60}
\end{align}
The air residual is
\begin{align}
\boxed{\begin{aligned}
r_{a,A}={}&C_{a,A}\dot{\widehat T}_{a,A}-\frac{T_o-\widehat T_{a,A}}{R_{Ao}}-\frac{\widehat T_{m,A}-\widehat T_{a,A}}{R_{Am}}\\
&-Q_{c,A}-(1-\eta_{r,A})Q_{r,A}.
\end{aligned}}
\tag{P2.61}
\end{align}
The mass residual is
\begin{align}
\boxed{\begin{aligned}
r_{m,A}={}&C_{m,A}\dot{\widehat T}_{m,A}-\frac{\widehat T_{a,A}-\widehat T_{m,A}}{R_{Am}}-\eta_{r,A}Q_{r,A}.
\end{aligned}}
\tag{P2.62}
\end{align}
Define
\begin{align}
\boxed{r_A=\begin{bmatrix}r_{a,A}\\r_{m,A}\end{bmatrix}.}
\tag{P2.63}
\end{align}
The physics loss is
\begin{align}
\boxed{\mathcal L_f=\frac{1}{N_{\rm tr}}\sum_{k\in\mathcal I_{\rm tr}}\left[\left(\frac{r_{a,A}(t_k)}{q_A^\star}\right)^2+\left(\frac{r_{m,A}(t_k)}{q_A^\star}\right)^2\right].}
\tag{P2.64}
\end{align}
The jointly estimated parameter set is
\begin{align}
\boxed{\theta_{A,2C}=\{R_{Ao},R_{Am},C_{a,A},C_{m,A},\eta_{r,A}\}.}
\tag{P2.65}
\end{align}
The hidden trajectory \(\widehat T_{m,A}(t)\) is a nuisance trajectory required for parameter identification, not a measured target.

\section*{16. Identity / IND 1R1C inverse PINNs}
IND requires two separate inverse PINNs:
\begin{align}
\boxed{\widehat T_D(t)=\mathcal N_{\phi_D}(t),}
\tag{P2.66}
\end{align}
and
\begin{align}
\boxed{\widehat T_K(t)=\mathcal N_{\phi_K}(t).}
\tag{P2.67}
\end{align}
Dining residual:
\begin{align}
\boxed{\begin{aligned}
r_D={}&C_D\dot{\widehat T}_D-\frac{T_o-\widehat T_D}{R_{Do}}-Q_{c,D}-\eta_{r,D}Q_{r,D}.
\end{aligned}}
\tag{P2.68}
\end{align}
Kitchen residual:
\begin{align}
\boxed{\begin{aligned}
r_K={}&C_K\dot{\widehat T}_K-\frac{T_o-\widehat T_K}{R_{Ko}}-Q_{c,K}-\eta_{r,K}Q_{r,K}.
\end{aligned}}
\tag{P2.69}
\end{align}
For Kitchen,
\begin{align}
\boxed{Q_{c,K}=Q_{ZIC,K}+Q_{AC,K},}
\tag{P2.70}
\end{align}
and
\begin{align}
\boxed{Q_{r,K}=Q_{ZIR,K}.}
\tag{P2.71}
\end{align}
The optimization problems are independent:
\begin{align}
\boxed{(\widehat\phi_D,\widehat\theta_D)=\arg\min\left[\lambda_y\mathcal L_{y,D}+\lambda_f\mathcal L_{f,D}\right],}
\tag{P2.72}
\end{align}
and
\begin{align}
\boxed{(\widehat\phi_K,\widehat\theta_K)=\arg\min\left[\lambda_y\mathcal L_{y,K}+\lambda_f\mathcal L_{f,K}\right].}
\tag{P2.73}
\end{align}
There is no \(R_{DK}\).

\section*{17. Identity / IND 2R2C inverse PINNs}
Dining network:
\begin{align}
\boxed{\widehat x_D=\begin{bmatrix}\widehat T_{a,D}\\\widehat T_{m,D}\end{bmatrix}.}
\tag{P2.74}
\end{align}
Kitchen network:
\begin{align}
\boxed{\widehat x_K=\begin{bmatrix}\widehat T_{a,K}\\\widehat T_{m,K}\end{bmatrix}.}
\tag{P2.75}
\end{align}
Dining air residual:
\begin{align}
\boxed{\begin{aligned}
r_{a,D}={}&C_{a,D}\dot{\widehat T}_{a,D}-\frac{T_o-\widehat T_{a,D}}{R_{Do}}-\frac{\widehat T_{m,D}-\widehat T_{a,D}}{R_{Dm}}\\
&-Q_{c,D}-(1-\eta_{r,D})Q_{r,D}.
\end{aligned}}
\tag{P2.76}
\end{align}
Dining mass residual:
\begin{align}
\boxed{\begin{aligned}
r_{m,D}={}&C_{m,D}\dot{\widehat T}_{m,D}-\frac{\widehat T_{a,D}-\widehat T_{m,D}}{R_{Dm}}-\eta_{r,D}Q_{r,D}.
\end{aligned}}
\tag{P2.77}
\end{align}
Kitchen air residual:
\begin{align}
\boxed{\begin{aligned}
r_{a,K}={}&C_{a,K}\dot{\widehat T}_{a,K}-\frac{T_o-\widehat T_{a,K}}{R_{Ko}}-\frac{\widehat T_{m,K}-\widehat T_{a,K}}{R_{Km}}\\
&-Q_{c,K}-(1-\eta_{r,K})Q_{r,K}.
\end{aligned}}
\tag{P2.78}
\end{align}
Kitchen mass residual:
\begin{align}
\boxed{\begin{aligned}
r_{m,K}={}&C_{m,K}\dot{\widehat T}_{m,K}-\frac{\widehat T_{a,K}-\widehat T_{m,K}}{R_{Km}}-\eta_{r,K}Q_{r,K}.
\end{aligned}}
\tag{P2.79}
\end{align}
Again, the two zone models are fitted independently.

\section*{18. Identity / DEP1 1R1C inverse PINN}
One joint PINN represents both measured temperatures:
\begin{align}
\boxed{\widehat x_\phi(t)=\begin{bmatrix}\widehat T_D(t)\\\widehat T_K(t)\end{bmatrix}.}
\tag{P2.80}
\end{align}
Dining residual:
\begin{align}
\boxed{\begin{aligned}
r_D={}&C_D\dot{\widehat T}_D-\frac{T_o-\widehat T_D}{R_{Do}}-\frac{\widehat T_K-\widehat T_D}{R_{DK}}-Q_{c,D}-\eta_{r,D}Q_{r,D}.
\end{aligned}}
\tag{P2.81}
\end{align}
Kitchen residual:
\begin{align}
\boxed{\begin{aligned}
r_K={}&C_K\dot{\widehat T}_K-\frac{T_o-\widehat T_K}{R_{Ko}}-\frac{\widehat T_D-\widehat T_K}{R_{DK}}-Q_{c,K}-\eta_{r,K}Q_{r,K}.
\end{aligned}}
\tag{P2.82}
\end{align}
The shared coupling parameter is
\begin{align}
\boxed{R_{DK}=R_{KD}>0.}
\tag{P2.83}
\end{align}
The residual vector is
\begin{align}
\boxed{r_{\rm DEP1}^{1C}=\begin{bmatrix}r_D\\r_K\end{bmatrix}.}
\tag{P2.84}
\end{align}
The data loss is
\begin{align}
\boxed{\mathcal L_y=\frac{1}{N_{\rm tr}}\sum_{k\in\mathcal I_{\rm tr}}\left[\left(\frac{\widehat T_D-T_D}{s_{T_D}}\right)^2+\left(\frac{\widehat T_K-T_K}{s_{T_K}}\right)^2\right].}
\tag{P2.85}
\end{align}
The physics loss is
\begin{align}
\boxed{\mathcal L_f=\frac{1}{N_{\rm tr}}\sum_{k\in\mathcal I_{\rm tr}}\left[\left(\frac{r_D}{q_D^\star}\right)^2+\left(\frac{r_K}{q_K^\star}\right)^2\right].}
\tag{P2.86}
\end{align}
The parameter vector is
\begin{align}
\boxed{\theta_{\rm DEP1}^{1C}=\{R_{Do},R_{Ko},R_{DK},C_D,C_K,\eta_{r,D},\eta_{r,K}\}.}
\tag{P2.87}
\end{align}

\section*{19. Identity / DEP1 2R2C inverse PINN}
The joint PINN output is
\begin{align}
\boxed{\widehat x_\phi(t)=\begin{bmatrix}\widehat T_{a,D}\\\widehat T_{m,D}\\\widehat T_{a,K}\\\widehat T_{m,K}\end{bmatrix}.}
\tag{P2.88}
\end{align}
The observed output is
\begin{align}
\boxed{\widehat y=\begin{bmatrix}\widehat T_{a,D}\\\widehat T_{a,K}\end{bmatrix}.}
\tag{P2.89}
\end{align}
Dining-air residual:
\begin{align}
\boxed{\begin{aligned}
r_{a,D}={}&C_{a,D}\dot{\widehat T}_{a,D}-\frac{T_o-\widehat T_{a,D}}{R_{Do}}-\frac{\widehat T_{m,D}-\widehat T_{a,D}}{R_{Dm}}\\
&-\frac{\widehat T_{a,K}-\widehat T_{a,D}}{R_{DK}}-Q_{c,D}-(1-\eta_{r,D})Q_{r,D}.
\end{aligned}}
\tag{P2.90}
\end{align}
Dining-mass residual:
\begin{align}
\boxed{\begin{aligned}
r_{m,D}={}&C_{m,D}\dot{\widehat T}_{m,D}-\frac{\widehat T_{a,D}-\widehat T_{m,D}}{R_{Dm}}-\eta_{r,D}Q_{r,D}.
\end{aligned}}
\tag{P2.91}
\end{align}
Kitchen-air residual:
\begin{align}
\boxed{\begin{aligned}
r_{a,K}={}&C_{a,K}\dot{\widehat T}_{a,K}-\frac{T_o-\widehat T_{a,K}}{R_{Ko}}-\frac{\widehat T_{m,K}-\widehat T_{a,K}}{R_{Km}}\\
&-\frac{\widehat T_{a,D}-\widehat T_{a,K}}{R_{DK}}-Q_{c,K}-(1-\eta_{r,K})Q_{r,K}.
\end{aligned}}
\tag{P2.92}
\end{align}
Kitchen-mass residual:
\begin{align}
\boxed{\begin{aligned}
r_{m,K}={}&C_{m,K}\dot{\widehat T}_{m,K}-\frac{\widehat T_{a,K}-\widehat T_{m,K}}{R_{Km}}-\eta_{r,K}Q_{r,K}.
\end{aligned}}
\tag{P2.93}
\end{align}
The residual vector is
\begin{align}
\boxed{r_{\rm DEP1}^{2C}=\begin{bmatrix}r_{a,D}\\r_{m,D}\\r_{a,K}\\r_{m,K}\end{bmatrix}.}
\tag{P2.94}
\end{align}
The physics loss is
\begin{align}
\boxed{\begin{aligned}
\mathcal L_f=\frac{1}{N_{\rm tr}}\sum_{k\in\mathcal I_{\rm tr}}\Bigg[&\frac{r_{a,D}^2+r_{m,D}^2}{(q_D^\star)^2}+\frac{r_{a,K}^2+r_{m,K}^2}{(q_K^\star)^2}\Bigg].
\end{aligned}}
\tag{P2.95}
\end{align}
The parameter vector is
\begin{align}
\boxed{\begin{aligned}
\theta_{\rm DEP1}^{2C}=\{&R_{Do},R_{Ko},R_{DK},R_{Dm},R_{Km},C_{a,D},C_{m,D},C_{a,K},C_{m,K},\\
&\eta_{r,D},\eta_{r,K}\}.
\end{aligned}}
\tag{P2.96}
\end{align}

\section*{20. Identity / DEP2 1R1C inverse PINN}
The neural trajectory is
\begin{align}
\boxed{\widehat x_\phi=\begin{bmatrix}\widehat T_D\\\widehat T_K\end{bmatrix}.}
\tag{P2.97}
\end{align}
Define
\begin{align}
\bar Q_c^{nh}&=Q_{ZIC,A}+Q_{Sol1,A},\\
\bar Q_r&=Q_{ZIR,A}+Q_{Sol2,A}.
\end{align}
Dining residual:
\begin{align}
\boxed{\begin{aligned}
r_D={}&C_D\dot{\widehat T}_D-\frac{T_o-\widehat T_D}{R_{Do}}-\frac{\widehat T_K-\widehat T_D}{R_{DK}}-Q_{AC,D}\\
&-\lambda_{c,D}\bar Q_c^{nh}-\eta_{r,D}\lambda_{r,D}\bar Q_r.
\end{aligned}}
\tag{P2.98}
\end{align}
Kitchen residual:
\begin{align}
\boxed{\begin{aligned}
r_K={}&C_K\dot{\widehat T}_K-\frac{T_o-\widehat T_K}{R_{Ko}}-\frac{\widehat T_D-\widehat T_K}{R_{DK}}-Q_{AC,K}\\
&-\lambda_{c,K}\bar Q_c^{nh}-\eta_{r,K}\lambda_{r,K}\bar Q_r.
\end{aligned}}
\tag{P2.99}
\end{align}
For the primary identification,
\begin{align}
\boxed{\eta_{r,D}=\eta_{r,K}=1,}
\tag{P2.100}
\end{align}
so \(\alpha_c,\alpha_r\) are learned without the \(\eta_r\lambda_r\) confounding. The optimization variables include
\begin{align}
\boxed{\Theta_{\rm train}=\{\phi,\rho_R,\rho_C,\alpha_c,\alpha_r\}.}
\tag{P2.101}
\end{align}

\section*{21. Identity / DEP2 2R2C inverse PINN}
State approximation:
\begin{align}
\boxed{\widehat x_\phi=\begin{bmatrix}\widehat T_{a,D}\\\widehat T_{m,D}\\\widehat T_{a,K}\\\widehat T_{m,K}\end{bmatrix}.}
\tag{P2.102}
\end{align}
Dining air:
\begin{align}
\boxed{\begin{aligned}
r_{a,D}={}&C_{a,D}\dot{\widehat T}_{a,D}-\frac{T_o-\widehat T_{a,D}}{R_{Do}}-\frac{\widehat T_{m,D}-\widehat T_{a,D}}{R_{Dm}}\\
&-\frac{\widehat T_{a,K}-\widehat T_{a,D}}{R_{DK}}-Q_{AC,D}-\lambda_{c,D}\bar Q_c^{nh}\\
&-(1-\eta_{r,D})\lambda_{r,D}\bar Q_r.
\end{aligned}}
\tag{P2.103}
\end{align}
Dining mass:
\begin{align}
\boxed{\begin{aligned}
r_{m,D}={}&C_{m,D}\dot{\widehat T}_{m,D}-\frac{\widehat T_{a,D}-\widehat T_{m,D}}{R_{Dm}}-\eta_{r,D}\lambda_{r,D}\bar Q_r.
\end{aligned}}
\tag{P2.104}
\end{align}
Kitchen air:
\begin{align}
\boxed{\begin{aligned}
r_{a,K}={}&C_{a,K}\dot{\widehat T}_{a,K}-\frac{T_o-\widehat T_{a,K}}{R_{Ko}}-\frac{\widehat T_{m,K}-\widehat T_{a,K}}{R_{Km}}\\
&-\frac{\widehat T_{a,D}-\widehat T_{a,K}}{R_{DK}}-Q_{AC,K}-\lambda_{c,K}\bar Q_c^{nh}\\
&-(1-\eta_{r,K})\lambda_{r,K}\bar Q_r.
\end{aligned}}
\tag{P2.105}
\end{align}
Kitchen mass:
\begin{align}
\boxed{\begin{aligned}
r_{m,K}={}&C_{m,K}\dot{\widehat T}_{m,K}-\frac{\widehat T_{a,K}-\widehat T_{m,K}}{R_{Km}}-\eta_{r,K}\lambda_{r,K}\bar Q_r.
\end{aligned}}
\tag{P2.106}
\end{align}
The residual vector is
\begin{align}
\boxed{r_{\rm DEP2}^{2C}=\begin{bmatrix}r_{a,D}\\r_{m,D}\\r_{a,K}\\r_{m,K}\end{bmatrix}.}
\tag{P2.107}
\end{align}

\section*{22. The complete inverse-PINN loss}
For every architecture,
\begin{align}
\boxed{\mathcal L_{\rm InvPINN}=\lambda_y\mathcal L_y+\lambda_f\mathcal L_f.}
\tag{P2.108}
\end{align}
Here \(\mathcal L_y\) contains only measured air-temperature mismatch, while \(\mathcal L_f\) contains the complete physical RC residual. There is no artificial loss on unmeasured mass temperature, no separate state decoder, no learned dynamics residual, and no projection.

\section*{23. Why no separate initial-condition loss is required during identification}
Observed air temperatures are already constrained at all training timestamps through \(\mathcal L_y\). For 2R2C, \(T_m\) is unobserved, so we do not impose
\begin{align}
T_m=T_a
\end{align}
as a physical identity. The PINN infers \(\widehat T_m(t)\) through the two RC residual equations. Initialization near
\begin{align}
T_m\approx T_a
\end{align}
may be used numerically, but it is not imposed as a model equation.

\section*{24. Collocation set}
The primary baseline uses
\begin{align}
\boxed{\mathcal C=\{t_k:k\in\mathcal I_{\rm tr}\}.}
\tag{P2.109}
\end{align}
Thus every training observation timestamp is also a physics collocation point. We do not introduce synthetic intermediate-time collocation points in the primary baseline, so no interpolation law must be invented for the measured forcing channels.

\section*{25. What is differentiated and what is not}
At each collocation point \(t_k\), \(\widehat x_\phi(t_k)\) is differentiated with respect to time. The forcing measurements \(v_k\) are not differentiated. In particular,
\begin{align}
Q_{AC}(t_k)
\end{align}
is treated as a known input to the balance, and we do not require
\begin{align}
\dot Q_{AC},\qquad \dot T_o,\qquad \dot Q_{ZIC},\qquad \dot Q_{ZIR}.
\end{align}

\section*{26. Training variables by architecture}
All-to-One 1R1C:
\begin{align}
\boxed{\phi,\;R_{Ao},\;C_A,\;[\eta_{r,A}]}
\tag{P2.110}
\end{align}
All-to-One 2R2C:
\begin{align}
\boxed{\phi,\;R_{Ao},\;R_{Am},\;C_{a,A},\;C_{m,A},\;[\eta_{r,A}].}
\tag{P2.111}
\end{align}
IND 1R1C Dining:
\begin{align}
\boxed{\phi_D,\;R_{Do},\;C_D,\;[\eta_{r,D}].}
\tag{P2.112}
\end{align}
IND 1R1C Kitchen:
\begin{align}
\boxed{\phi_K,\;R_{Ko},\;C_K,\;[\eta_{r,K}].}
\tag{P2.113}
\end{align}
IND 2R2C per zone \(i\):
\begin{align}
\boxed{\phi_i,\;R_{io},\;R_{im},\;C_{a,i},\;C_{m,i},\;[\eta_{r,i}].}
\tag{P2.114}
\end{align}
DEP1 1R1C:
\begin{align}
\boxed{\phi,\;R_{Do},R_{Ko},R_{DK},C_D,C_K,[\eta_{r,D},\eta_{r,K}].}
\tag{P2.115}
\end{align}
DEP1 2R2C:
\begin{align}
\boxed{\begin{aligned}
&\phi,R_{Do},R_{Ko},R_{DK},R_{Dm},R_{Km},\\
&C_{a,D},C_{m,D},C_{a,K},C_{m,K},[\eta_{r,D},\eta_{r,K}].
\end{aligned}}
\tag{P2.116}
\end{align}
DEP2 adds
\begin{align}
\boxed{\alpha_c,\alpha_r.}
\tag{P2.117}
\end{align}

\section*{27. Training procedure}
\subsection*{Stage A --- trajectory prefit}
Optimize only \(\phi\):
\begin{align}
\boxed{\min_\phi\mathcal L_y.}
\tag{P2.118}
\end{align}
For 2R2C, air outputs are fitted to measured temperatures; hidden mass outputs are initialized numerically near the corresponding air temperatures but receive no artificial labels.

\subsection*{Stage B --- joint inverse-PINN optimization}
Then jointly optimize neural and physical variables:
\begin{align}
\boxed{\min_{\phi,\theta}\lambda_y\mathcal L_y+\lambda_f\mathcal L_f.}
\tag{P2.119}
\end{align}
DEP2 also optimizes \(\alpha_c,\alpha_r\).

\subsection*{Stage C --- parameter-focused refinement}
Optionally freeze \(\phi\) and refine only physical parameters:
\begin{align}
\boxed{\min_\theta\mathcal L_f\left(\widehat x_{\widehat\phi},\dot{\widehat x}_{\widehat\phi};\theta\right).}
\tag{P2.120}
\end{align}
This is a numerical refinement, not a different model.

\section*{28. Validation must use the physical RC model, not the PINN trajectory network}
Suppose epoch \(e\) gives parameters \(\theta^{(e)}\). Construct
\begin{align}
\boxed{\dot x=f_{RC}\left(x,v;\theta^{(e)}\right).}
\tag{P2.121}
\end{align}
Integrate that physical ODE over validation trajectories and score \(Hx_{RC}(t)\) against validation air temperature. Do not select checkpoints based on \(H\widehat x_\phi(t)\) at validation times. The selected parameter set is
\begin{align}
\boxed{\widehat\theta=\theta^{(e^\star)},}
\tag{P2.122}
\end{align}
where
\begin{align}
\boxed{e^\star=\arg\min_e\mathcal E_{\rm val}^{RC}\left(\theta^{(e)}\right).}
\tag{P2.123}
\end{align}

\section*{29. 1R1C validation initialization}
For 1R1C the entire state is measured, so
\begin{align}
\boxed{x(t_0)=y(t_0).}
\tag{P2.124}
\end{align}

\section*{30. 2R2C validation/test initialization}
For 2R2C, \(T_a(t_0)\) is known but \(T_m(t_0)\) is hidden. Use a past warm-up window
\begin{align}
\mathcal W_0=\{t_{-L_w},\ldots,t_0\}.
\end{align}
For one zone,
\begin{align}
x_{-L_w}(\zeta)=\begin{bmatrix}T_a(t_{-L_w})\\\zeta\end{bmatrix}.
\end{align}
Let \(\Phi_\theta\) denote numerical integration of the identified RC model. Estimate
\begin{align}
\boxed{\widehat\zeta=\arg\min_\zeta\sum_{\ell=-L_w}^{0}\left\|H\Phi_{\widehat\theta}\left(t_\ell;t_{-L_w},x_{-L_w}(\zeta),v\right)-y(t_\ell)\right\|_2^2.}
\tag{P2.125}
\end{align}
Then
\begin{align}
\boxed{\widehat T_m(t_0)=\left[\Phi_{\widehat\theta}\left(t_0;t_{-L_w},x_{-L_w}(\widehat\zeta),v\right)\right]_m.}
\tag{P2.126}
\end{align}
For two-zone 2R2C,
\begin{align}
\zeta=\begin{bmatrix}T_{m,D}(t_{-L_w})\\T_{m,K}(t_{-L_w})\end{bmatrix}.
\end{align}
If no warm-up history exists, the neutral fallback is
\begin{align}
\boxed{T_m(t_0)=T_a(t_0).}
\tag{P2.127}
\end{align}
This fallback is an initialization convention, not an RC equation.

\section*{31. Final deployed model after inverse-PINN training}
After selecting \(\widehat\theta\), the PINN parameters \(\widehat\phi\) are discarded. For 1R1C,
\begin{align}
\boxed{\dot x=f_{RC}^{1C}\left(x,v;\widehat\theta\right).}
\tag{P2.128}
\end{align}
For 2R2C,
\begin{align}
\boxed{\dot x=f_{RC}^{2C}\left(x,v;\widehat\theta\right).}
\tag{P2.129}
\end{align}
The deployed model contains no neural vector field.

\section*{32. Sim1 meaning for inverse PINN-RC}
For 1R1C, \(x_k=y_k\). For 2R2C, the observed air state is measured and the hidden mass state comes from the causal initialization/propagation mechanism above. With measured forcing \(v_k\),
\begin{align}
\boxed{\widehat x_{k+1}=\Phi_{\widehat\theta}\left(x_k,v_k,\Delta t\right).}
\tag{P2.130}
\end{align}
and
\begin{align}
\boxed{\widehat y_{k+1}=H\widehat x_{k+1}.}
\tag{P2.131}
\end{align}

\section*{33. Sim2 meaning for inverse PINN-RC}
Initialize once and recursively compute
\begin{align}
\boxed{\widehat x_{k+1}=\Phi_{\widehat\theta}\left(\widehat x_k,v_k,\Delta t\right).}
\tag{P2.132}
\end{align}
The model's own predicted state is fed back; recorded Phase-D QAC and disturbances remain external forcing.

\section*{34. Sim3 meaning for inverse PINN-RC}
The thermostat produces \(\dot m_k\). Phase C produces
\begin{align}
\boxed{Q_{AC,k}=G_{QAC}\left(\dot m_k,\widehat T_{z,k},T_{sa,k},\ldots\right).}
\tag{P2.133}
\end{align}
Then
\begin{align}
\boxed{\widehat x_{k+1}=\Phi_{\widehat\theta}\left(\widehat x_k,Q_{AC,k},d_k\right).}
\tag{P2.134}
\end{align}
The inverse PINN itself is not involved in Sim3 after parameter identification.

\section*{35. What we log during inverse-PINN training}
Record
\begin{align}
\boxed{\mathcal L_y},\qquad \boxed{\mathcal L_f},\qquad \boxed{\mathcal L_{\rm InvPINN}}.
\end{align}
For every trainable physical parameter, record by epoch
\begin{align}
R_j^{(e)},\qquad C_j^{(e)},\qquad \eta_j^{(e)}.
\end{align}
For DEP2 also record
\begin{align}
\lambda_{c,D}^{(e)},\lambda_{c,K}^{(e)},\lambda_{r,D}^{(e)},\lambda_{r,K}^{(e)}.
\end{align}
Log residual RMSE separately for every physical equation, e.g.
\begin{align}
\operatorname{RMSE}(r_{a,D}),\quad \operatorname{RMSE}(r_{m,D}),\quad \operatorname{RMSE}(r_{a,K}),\quad \operatorname{RMSE}(r_{m,K}).
\end{align}

\section*{36. No identifiability claim}
We do not claim
\begin{align}
\boxed{\widehat\theta=\theta_{\rm true}}
\end{align}
in any universal sense. Instead,
\begin{align}
\boxed{\widehat\theta\text{ is the physically constrained parameter set identified from the fixed training data and architecture.}}
\tag{P2.135}
\end{align}
The paper evaluates whether the resulting identified RC model predicts and controls well.

\section*{37. Multi-start estimation}
For restart
\begin{align}
s=1,\ldots,S,
\end{align}
obtain \(\widehat\theta^{(s)}\) and select using physical-model validation error:
\begin{align}
\boxed{s^\star=\arg\min_s\mathcal E_{\rm val}^{RC}\left(\widehat\theta^{(s)}\right).}
\tag{P2.136}
\end{align}
Then
\begin{align}
\boxed{\widehat\theta=\widehat\theta^{(s^\star)}.}
\tag{P2.137}
\end{align}
The test set remains untouched until selection is complete.

\section*{38. Exact distinction from the other three methods}
Inverse PINN-RC approximates
\begin{align}
\boxed{x(t)}
\end{align}
during identification and retains the final vector field
\begin{align}
\boxed{f_{RC}(x,v;\widehat\theta).}
\end{align}
NODE will approximate
\begin{align}
\boxed{f_\omega(x,v)}
\end{align}
directly. Base PINODE will also approximate \(f_\omega(x,v)\), but with a soft RC residual penalty. EBP-PINODE will approximate a raw derivative
\begin{align}
\widetilde f_\omega(x,v)
\end{align}
which is then physically projected.

\section*{39. Complete inverse-PINN algorithm}
Choose
\begin{align}
\text{architecture}\in\{A,\mathrm{IND},\mathrm{DEP1},\mathrm{DEP2}\}
\end{align}
and
\begin{align}
\text{RC order}\in\{1R1C,2R2C\}.
\end{align}
Load only \(\mathcal I_{\rm tr}\), compute training-only normalization statistics, construct \(\widehat x_\phi(t)\), obtain \(\dot{\widehat x}_\phi(t)\) by automatic differentiation, and transform unconstrained variables through
\begin{align}
\rho_R&\rightarrow R>0,\\
\rho_C&\rightarrow C>0,\\
\rho_\eta&\rightarrow\eta\in(0,1).
\end{align}
For DEP2,
\begin{align}
\alpha_c,\alpha_r\rightarrow\lambda_c,\lambda_r.
\end{align}
Evaluate the exact architecture-specific RC residuals and minimize
\begin{align}
\boxed{\mathcal L_{\rm InvPINN}=\lambda_y\mathcal L_y+\lambda_f\mathcal L_f.}
\tag{P2.138}
\end{align}
At candidate checkpoints, validate the physical RC ODE directly, choose the best physical parameter checkpoint, discard \(\mathcal N_\phi\), freeze \(\widehat\theta\), and use only the identified RC ODE for final tests.

\section*{Part 2 conclusion}
The inverse PINN baseline is not
\begin{align}
x_{k+1}=\text{NN}(x_k,u_k,d_k).
\end{align}
It is not a Neural ODE and not a hybrid RC+NN deployment model. It is an inverse physical-parameter identification method:
\begin{align}
\boxed{\text{PINN trajectory approximation}+\text{exact RC residual}\longrightarrow\widehat R,\widehat C,\widehat\eta_r,[\widehat\lambda]}
\end{align}
followed by
\begin{align}
\boxed{\text{PINN discarded}}
\end{align}
and
\begin{align}
\boxed{\dot x=f_{RC}(x,v;\widehat\theta)}
\end{align}
used as the final interpretable simulation/control model.

The DEP2 1R1C safeguard is
\begin{align}
\boxed{\text{do not simultaneously learn }\eta_r\text{ and }\lambda_r\text{ in the primary DEP2 1R1C fit}.}
\end{align}
For the primary DEP2 1R1C experiment,
\begin{align}
\boxed{\eta_r=1,\qquad\lambda_r\ \text{learnable}.}
\end{align}
Everything else follows directly from the Part-1 architecture.

\end{document}
